% TOP_PROBLEM: {"id": 30006886, "title": "Fujita Freeness Conjecture", "category_id": 5, "category": {"id": 5, "name": "algebraic_geometry", "display_name": "Algebraic Geometry", "description": "Geometric objects defined by polynomial equations.", "slug": "algebraic-geometry", "order_index": 5, "created_at": "2026-07-31T15:26:25.671Z"}, "status": "open"}
% FIRST_POSTED: 2026-09-25
% ENTRY_KIND: statement_only
% !TeX program = lualatex
% Definition and sources only; this file is not an LLM solution attempt.
\documentclass[11pt]{article}
\usepackage[margin=25mm]{geometry}
\usepackage{fontspec}
\setmainfont{Latin Modern Roman}
\usepackage{amsmath,amssymb,mathtools}
\usepackage{unicode-math}
\setmathfont{Latin Modern Math}
\usepackage{xurl,hyperref}
\hypersetup{colorlinks=true,urlcolor=blue}
\setcounter{secnumdepth}{0}
\setlength{\parindent}{0pt}
\setlength{\parskip}{5pt}
\setlength{\emergencystretch}{3em}
\begin{document}
\section*{\texorpdfstring{Fujita Freeness Conjecture}{Fujita Freeness Conjecture}}\label{30006886}
\subsection{Definitions and mathematical statement}
Let $X$ be a smooth complex projective variety of dimension $n$, $K_X=\bigwedge^n\Omega_X^1$ its canonical line bundle, and $L$ an ample line bundle. Ampleness means that some positive tensor power defines a projective embedding. A line bundle $A$ is basepoint-free if the evaluation map $H^0(X,A)\otimes\mathcal O_X\to A$ is surjective, equivalently every point has a global section nonvanishing there. Fujita freeness asserts
\[
 K_X\otimes L^{\otimes(n+1)}\text{ is basepoint-free}.
\]
In additive divisor notation this is $K_X+(n+1)L$. The assertion requires no extra positivity such as very ampleness of $L$, and it does not claim that this adjoint bundle itself gives an embedding; the very-ampleness target is listed separately.
\par\smallskip\textit{Formulation and concept references:} \href{https://arxiv.org/abs/2408.06530v2}{[S1]}.

\subsection{Short English statement}
Does adding $(n+1)$ copies of any ample divisor to the canonical divisor of a smooth projective $n$-fold always remove every base point?
\subsection{Sources}
\begin{itemize}
\item[S1]Takumi Murayama, Moving Seshadri constants and effective Fujita-type conjectures, arXiv:2408.06530v2, Conjecture 1.1.1(i).
\url{https://arxiv.org/abs/2408.06530v2}.
\textit{Formulation source.}
\end{itemize}
\end{document}
